% Jade Zhao — how I think about mentorship, clients, and shipping usable systems.

\documentclass[11pt]{article}

\usepackage[margin=1in]{geometry}
\usepackage{hyperref}
\usepackage{parskip}

\pagestyle{empty}
\hypersetup{colorlinks=true,urlcolor=blue}

\begin{document}

% Shared header fragment — \input{} inside \begin{document} of each standalone Jade Zhao document.
% Contact block matches images.tex / site; omit phone (not in repo).

\begin{center}
    {\LARGE \textbf{Jade Zhao}}\\
    \vspace{0.25em}
    {\small Bloomington, IN \;|\; Madrid, Spain (Spring 2026)}\\
    \vspace{0.25em}
    {\small
      \href{mailto:jlzhao@iu.edu}{jlzhao@iu.edu}
      \;|\;
      \href{https://linkedin.com/in/jadexzhao}{linkedin.com/in/jadexzhao}
      \;|\;
      \href{https://github.com/jadewowgreen}{github.com/jadewowgreen}
    }
\end{center}


\vspace{0.5em}
\noindent\rule{\linewidth}{1pt}
\vspace{0.75em}

\noindent\today

\vspace{0.75em}

I treat student-facing and client-facing technical work as one craft: respect people's time, tell the truth about scope, and leave them with a clear next step. Whether I am answering mentee questions through IU's mentoring programs, leading a ServeIT Agile team, or wiring database workflows for health informatics research, the bar is the same---can this person trust what we shipped tomorrow morning?

I care about educational and civic contexts where software reduces administrative drag instead of adding it. That has meant consolidating campus resources into accessible tools, automating repetitive advisor communication, and helping students navigate international applications with better structure than a PDF graveyard.

Good relationships go both ways. I bring documentation, follow-through, and honest estimates; students and stakeholders bring goals and pushback that should change the roadmap. I translate jargon into something you can act on so ``I understood the meeting'' actually becomes ``I know what to do on Monday.''

In research and product-shaped projects I default to agency: interfaces and data flows where humans can correct mistakes, and AI features paired with bias checks, audit trails, and educator oversight so ``innovation'' does not mean ``no one is accountable.''

Health informatics, nonprofit redesigns, and clinic work have all forced me to align engineers, domain experts, and end users under shared definitions of risk---there is no solo hero version of that job I believe in.

Accessibility and first-gen student support belong in the same backlog as every other feature. Audits, WCAG fixes, and roadmaps are how I make access real instead of rhetorical.

\vspace{1.5em}
\noindent\textit{Best regards,}\\
\textbf{Jade Zhao}

\end{document}
